\chapter{Binge Eating Disorder}

\section*{Definition and Overview}
Binge eating disorder (BED) is a serious eating disorder in which a person regularly consumes an unusually large amount of food within a discrete period while experiencing a sense of loss of control over their eating. Unlike bulimia nervosa, binge episodes in BED are not followed by purging, vomiting, fasting, or excessive exercise.

People with BED often eat much more rapidly than normal, continue eating until uncomfortably full, and feel intense shame, guilt, or disgust afterwards. It is the most common eating disorder, affects people of all body sizes, and is a treatable condition for which early professional help substantially improves the outlook.

\section*{Epidemiology}
Binge eating disorder is more common than is widely appreciated and affects people of all ages, sexes, and backgrounds.

\begin{itemize}
    \item \textbf{Lifetime prevalence:} roughly 2 to 3 per cent of people will experience BED at some point in their lives
    \item \textbf{Sex distribution:} it is somewhat more common in women, but men are significantly affected
    \item \textbf{Age at onset:} the disorder usually begins in late adolescence or early adulthood
    \item \textbf{Weight status:} BED is strongly associated with overweight and obesity, although people of normal weight can also be affected
    \item \textbf{Co-morbidity:} depression, anxiety, and other mental health conditions frequently co-occur
    \item \textbf{Health-seeking:} many people never seek help, partly because of shame and stigma
\end{itemize}

\section*{Aetiology and Pathophysiology}
Binge eating disorder arises from an interaction of genetic, biological, psychological, and social factors.

\begin{itemize}
    \item \textbf{Genetics:} the condition runs in families, and twin studies indicate substantial heritability
    \item \textbf{Brain chemistry:} differences in the brain's reward and appetite pathways, including dopamine signalling, may make binge eating more rewarding and harder to resist
    \item \textbf{Psychological factors:} low mood, anxiety, poor body image, and perfectionism are common among people with BED
    \item \textbf{Dieting:} restrictive dieting and rigid food rules often trigger binge episodes, and binges frequently occur during or after diet attempts
    \item \textbf{Social and cultural pressures:} weight stigma and sociocultural pressure to be thin contribute to distress about eating
    \item \textbf{Life events:} stress, trauma, and difficult emotions frequently precipitate individual binge episodes
\end{itemize}

\section*{Clinical Features}
The core feature is recurrent binge eating, defined as eating a large amount of food in a discrete period with a sense of loss of control. Other features include:

\begin{itemize}
    \item Eating much more rapidly than normal during a binge episode
    \item Eating until uncomfortably full
    \item Eating large amounts of food when not physically hungry
    \item Eating alone or hiding food because of embarrassment about the amount consumed
    \item Feeling disgusted with oneself, depressed, or very guilty after a binge
    \item The absence of regular compensatory behaviours, such as self-induced vomiting or misuse of laxatives
    \item Significant distress about the eating behaviour
    \item Physical effects of excess weight, including fatigue, joint pain, poor sleep, and breathlessness on exertion
\end{itemize}

\section*{Differential Diagnosis}
Conditions that can resemble binge eating disorder include:

\begin{itemize}
    \item \textbf{Bulimia nervosa:} binge eating followed by compensatory behaviours such as vomiting, fasting, or misuse of laxatives
    \item \textbf{Depression:} some people with depression overeat, but usually without discrete binge episodes and the accompanying loss of control
    \item \textbf{Emotional or reactive overeating:} eating more than intended at times of stress, without the large, uncontrolled episodes that define BED
    \item \textbf{Obesity without an eating disorder:} weight gain alone does not mean BED is present
    \item \textbf{Medication side effects:} some medicines, including certain antidepressants and antipsychotics, increase appetite and weight
    \item \textbf{Night eating syndrome and other specified feeding or eating disorders:} related but distinct conditions
\end{itemize}

\section*{When to Seek Medical Care}
See a doctor if you or someone you care about regularly loses control when eating, feels distressed about eating, or has binge episodes at least weekly. Early help greatly improves the chance of recovery.

\textbf{Call 000 (or local emergency services) for:}
\begin{itemize}
    \item Thoughts of suicide or self-harm, or plans to act on them
    \item Chest pain, collapse, or fainting
    \item Severe breathlessness or palpitations
    \item Any other sudden, severe physical or mental health emergency
\end{itemize}

In addition to emergency services, free 24-hour crisis support is available by telephone from services such as Lifeline (13 11 14).

\section*{Diagnosis and Investigations}
\begin{itemize}
    \item \textbf{Clinical interview:} a careful, non-judgmental discussion of eating patterns, binge episodes, and associated feelings of loss of control and distress
    \item \textbf{Assessment against diagnostic criteria:} binge episodes with a sense of loss of control occurring at least once a week for three months, without compensatory behaviours
    \item \textbf{Physical examination:} measurement of weight and height, blood pressure, and a general health assessment
    \item \textbf{Blood tests:} blood sugar, cholesterol, and liver function, because BED is associated with metabolic conditions such as type 2 diabetes and fatty liver disease
    \item \textbf{Screening for related conditions:} depression, anxiety, and other eating disorders are common and require separate assessment
    \item \textbf{Weight history:} review of past dieting, weight fluctuations, and any physical complications
\end{itemize}

\section*{Management}
Binge eating disorder is treatable, and most people improve significantly with professional help.

\begin{itemize}
    \item \textbf{Cognitive behaviour therapy (CBT-E):} the first-line treatment, which helps people break the binge cycle by regularising eating patterns and addressing the thoughts that drive binges
    \item \textbf{Guided self-help:} a structured cognitive behaviour therapy programme with brief professional support, which can be effective, particularly in the early stages
    \item \textbf{Interpersonal psychotherapy:} an alternative therapy focusing on the relationship and mood issues that can drive binge eating
    \item \textbf{Medication:} lisdexamfetamine is approved in some countries for moderate to severe BED; antidepressant medication may be used for co-occurring depression
    \item \textbf{Dietitian support:} advice on regular, balanced meals and on sustainable weight management, avoiding restrictive diets that trigger binges
    \item \textbf{Treatment of co-morbid conditions:} management of depression, anxiety, diabetes, and sleep apnoea where present
\end{itemize}

\section*{Complications}
\begin{itemize}
    \item \textbf{Obesity and its consequences:} type 2 diabetes, high blood pressure, heart disease, obstructive sleep apnoea, and joint problems
    \item \textbf{Metabolic syndrome and fatty liver disease}
    \item \textbf{Digestive problems:} bloating, indigestion, gall bladder disease, and discomfort from very large meals
    \item \textbf{Mental health effects:} depression, anxiety, low self-esteem, and worsening body image
    \item \textbf{Social and occupational effects:} isolation, relationship difficulties, and impaired performance at work or study
    \item \textbf{Mortality:} an increased risk of early death compared with the general population, largely from obesity-related disease
\end{itemize}

\section*{Prognosis and Outcomes}
The outlook for binge eating disorder is good with treatment.

\begin{itemize}
    \item Many people achieve significant reductions in binge episodes with cognitive behaviour therapy
    \item Recovery is often gradual, and setbacks can occur, but long-term improvement is common
    \item Untreated BED tends to persist for many years and is associated with worsening metabolic health
    \item Early treatment, regular eating patterns, and support for coexisting mood problems are associated with better outcomes
\end{itemize}

\section*{Prevention and Self-Care}
\begin{itemize}
    \item Aim for regular meals and snacks rather than skipping meals or extreme dieting, both of which trigger binges
    \item Eat mindfully and without guilt, and allow yourself all foods in moderation
    \item Develop healthy ways of coping with stress, such as physical activity, talking to others, or hobbies, rather than turning to food
    \item Challenge weight stigma and perfectionism, and focus on health rather than weight alone
    \item If binge episodes are beginning, seek help early rather than waiting for the pattern to become entrenched
    \item Parents and carers can help by modelling healthy attitudes to food and body image in the home
\end{itemize}

\section*{Sources and Further Reading}
The following organisations provide reliable information on binge eating disorder for patients, families, and health professionals.

\begin{itemize}
    \item NHS. \textit{Information on binge eating disorder, its symptoms, and treatment.} https://www.nhs.uk/
    \item Mayo Clinic. \textit{Overview of binge eating disorder and its management.} https://www.mayoclinic.org/
    \item NICE. \textit{Clinical guidance on the treatment of eating disorders, including binge eating disorder.} https://www.nice.org.uk/
    \item MedlinePlus. \textit{Health information on binge eating disorder for patients and families.} https://medlineplus.gov/
    \item National Institute of Mental Health. \textit{Educational resources on eating disorders.} https://www.nimh.nih.gov/
\end{itemize}
